\documentclass[11pt]{article}
\usepackage[margin=0.92in]{geometry}
\usepackage{amsmath,amssymb,amsthm}
\usepackage{graphicx}
\usepackage{microtype}
\usepackage{xcolor}
\usepackage[hidelinks]{hyperref}
\usepackage{enumitem}

\definecolor{deepblue}{RGB}{27,72,112}
\definecolor{softblue}{RGB}{235,243,249}
\newtheorem{theorem}{Theorem}
\newtheorem{lemma}{Lemma}
\newcommand{\R}{\mathbb R}
\newcommand{\Hh}{\mathcal H}
\newcommand{\dd}{\,d}
\newcommand{\cl}{\overline}

\title{Compact-Group Quotients Preserve Hilbert Representability\\
and an Affirmative Answer to Problem 8.1(c)}
\author{Candidate full solution -- run \texttt{fa\_banach\_001}, lane 14}
\date{9 August 2026}

\begin{document}
\maketitle

\begin{center}
\fcolorbox{deepblue}{softblue}{\parbox{0.89\textwidth}{
\textbf{Candidate full answer to option (c); likely valid.}
Every metrizable recurrent-transitive Hilbert system is an almost one-to-one
factor of a Hilbert-representable system.  Options (a), (b), and Problem 8.2
are not claimed.
}}
\end{center}

\section*{Source question}

Glasner and Weiss \cite{GW} prove in their Theorem 3.1 that every metrizable
recurrent-transitive Hilbert system has an almost one-to-one extension which
is a compact group-factor of a Hilbert-representable system.  On source PDF
page 13, Problem 8.1 asks whether either extension can be removed and lists
three precise alternatives.

\begin{figure}[h]
\centering
\includegraphics[width=\textwidth]{figures/open_problem_crop.png}
\caption{The structure-theorem context and Problems 8.1--8.2 in \cite{GW},
source PDF page 13.}
\end{figure}

\noindent\textbf{Selected transcription (Problem 8.1(c)).}
Is a metrizable recurrent-transitive Hilbert system always ``an almost 1-1
factor of a Hilbert-representable system''?

\section*{Definitions and source input}

A compact metric dynamical system $(X,T)$ is \emph{Hilbert-representable} if
there are a Hilbert space $\Hh$, a unitary operator $U$ on $\Hh$, and a weakly
compact $U$-invariant set $Z\subset\Hh$ such that $(X,T)$ is isomorphic to
$(Z,U|_Z)$.  A factor map is \emph{almost one-to-one} if the points in the
target having singleton fibers form a dense set (equivalently in this compact
metric setting, a dense $G_\delta$ set).

The only dynamical input used below is the construction in the proof of
\cite[Theorem 3.1]{GW}.  It gives a Hilbert-representable point-transitive
system $(X,U)$, a compact subgroup $K$ of its uniform Ellis group
$\Lambda_X$, and a diagram
\[
 X \longrightarrow \widetilde Y=X/K
 \mathop{\longrightarrow}^{\rho} Y,
\]
where the first map is the $K$-orbit quotient and $\rho$ is almost one-to-one.

\section*{Proof intuition}

The apparent loss of Hilbert representability occurs at the quotient $X/K$.
To recover it, replace a point $x$ by the probability measure uniformly
distributed on its compact orbit $Kx$.  Encode this measure by all of its
symmetric tensor moments.  The moments live in Hilbert tensor powers and
transform unitarily under the dynamics.  They separate orbit measures because
polynomials in weak linear coordinates are dense in the continuous functions
on the weakly compact set.  A weighted direct sum of all moments therefore
embeds $X/K$ into one weakly compact Hilbert system.

The compact group in the source is initially a group of homeomorphisms of
$X$.  A short cyclicity argument lifts it to strongly continuous commuting
unitaries, so the moment construction applies to exactly the quotient in
Theorem 3.1.

\section*{Compact-group quotient lemma}

Write $\Hh_{\R}$ for the realification of a real or complex Hilbert space, and
$\Hh_{\R}^{\odot n}$ for its $n$th real symmetric Hilbert tensor power.

\begin{lemma}[Tensor-moment quotient]
Let $X\subset\Hh$ be weakly compact and invariant under a unitary $U$.  Let
$K$ be a compact group with a strongly continuous unitary representation on
$\Hh$ such that every $k\in K$ preserves $X$ and commutes with $U$.  Then the
quotient dynamical system $(X/K,\bar U)$ is Hilbert-representable.
\end{lemma}

\begin{proof}
Pass to $\Hh_{\R}$, on which all the operators are orthogonal and the weak
topology on $X$ is unchanged.  Let $m_K$ be normalized Haar measure.  For
$n\geq1$ define
\[
 M_n(x)=\int_K (kx)^{\otimes n}\dd m_K(k)
 \quad\in\Hh_{\R}^{\odot n}.
\]
This Bochner integral exists because the orbit map is norm-continuous.  It is
constant on $K$-orbits by invariance of Haar measure.

Each $M_n:X\to\Hh_{\R}^{\odot n}$ is weak-to-weak continuous.  Indeed, for a
pure symmetric tensor $v_1\odot\cdots\odot v_n$, its scalar pairing with
$M_n(x)$ is an integral of
\[
 \prod_{j=1}^n \langle kx,v_j\rangle
 =\prod_{j=1}^n \langle x,k^{-1}v_j\rangle.
\]
This is jointly continuous in $(k,x)\in K\times X$: the vectors $k^{-1}v_j$
vary in norm and $X$ is bounded in norm.  Integration over compact $K$
preserves continuity.  Finite linear combinations of pure tensors are dense,
and the estimate $\|M_n(x)\|\leq R^n$, where
$R=\sup_{x\in X}\|x\|$, extends the conclusion to every test tensor.

The moments separate $K$-orbits.  To see this, let $\mu_x$ be the pushforward
of $m_K$ under $k\mapsto kx$.  If $M_n(x)=M_n(y)$ for all $n\geq1$, then
$\mu_x$ and $\mu_y$ have the same integrals against every finite product of
real weak linear coordinates.  The algebra of such real polynomials contains
the constants and separates points of the weakly compact Hausdorff space
$X$.  The real Stone--Weierstrass theorem \cite{Rudin} gives
$\mu_x=\mu_y$.  Haar measure has full support, so
\[
 \operatorname{supp}\mu_x=Kx,\qquad
 \operatorname{supp}\mu_y=Ky.
\]
Thus $Kx=Ky$.

Choose positive weights
$a_n=2^{-n}(1+R)^{-n}$ and set
\[
 \Phi(x)=\bigoplus_{n=1}^{\infty}a_nM_n(x)
 \quad\in\mathcal K:=\bigoplus_{n=1}^{\infty}\Hh_{\R}^{\odot n}.
\]
The bound $a_n\|M_n(x)\|\leq2^{-n}$ gives a uniformly square-summable tail,
so coordinatewise weak continuity implies weak continuity of $\Phi$.  The map
descends to a continuous injection $\bar\Phi:X/K\to\mathcal K$.  Since $X/K$
is compact and the weak topology is Hausdorff, this is a homeomorphism onto a
weakly compact set.

Finally, commutation of $U$ with $K$ gives
\[
 M_n(Ux)=U^{\odot n}M_n(x).
\]
Hence $\bar\Phi$ intertwines $\bar U$ with the orthogonal operator
$\bigoplus_{n\geq1}U^{\odot n}$ on $\mathcal K$.  This is a Hilbert
representation of the quotient system.
\end{proof}

\section*{Lifting the compact kernel}

\begin{lemma}[Uniform Ellis group lift]
Let $(X,U)$ be a point-transitive Hilbert-representable metric system.  After
choosing a cyclic Hilbert realization
\[
 X=\operatorname{weak\!\!-cl}\{U^nx_0:n\in\mathbb Z\},
 \qquad
 \Hh=\cl{\operatorname{span}}\{U^nx_0:n\in\mathbb Z\},
\]
the uniform Ellis group $\Lambda_X$ has a strongly continuous unitary
representation on $\Hh$.  It extends its action on $X$, and its image commutes
with $U$.
\end{lemma}

\begin{proof}
Every transitive point $z\in X$ has $\|z\|=\|x_0\|$.  Indeed, as a weak limit
of orbit points it satisfies $\|z\|\leq\|x_0\|$.  Since $x_0$ is in the weak
orbit closure of the transitive point $z$, weak lower semicontinuity gives the
reverse inequality.

For $g\in\Lambda_X$, choose integers $n_j$ for which $U^{n_j}|_X$ converges
uniformly to $g$ in a compatible weak metric.  Then
$U^{n_j}x_0\to gx_0$ weakly.  The point $gx_0$ is transitive, so the preceding
norm equality upgrades this convergence to norm convergence.  On every
vector $p(U)x_0$ in the algebraic cyclic span,
\[
 U^{n_j}p(U)x_0=p(U)U^{n_j}x_0
\]
therefore converges in norm.  Uniform boundedness and density give a strong
operator limit $\widehat g$ on $\Hh$.  Applying the same construction to
$g^{-1}$ shows that $\widehat g$ is unitary.  It is independent of the
approximating sequence, extends $g$ on $X$, and commutes with $U$.

If $g_i\to g$ in $\Lambda_X$, then $g_ix_0\to gx_0$ weakly; all these points
are transitive and have the same norm.  Thus the convergence is in norm, and
cyclicity gives $\widehat g_i\to\widehat g$ strongly.  The lift is strongly
continuous.
\end{proof}

\section*{Result}

\begin{theorem}[Affirmative answer to Problem 8.1(c)]
Every metrizable recurrent-transitive Hilbert system is an almost one-to-one
factor of a Hilbert-representable system.
\end{theorem}

\begin{proof}
Let $(Y,T)$ be such a system.  Apply the construction of
\cite[Theorem 3.1]{GW}.  It supplies a point-transitive
Hilbert-representable system $(X,U)$, a compact subgroup
$K\leq\Lambda_X$, the quotient $\widetilde Y=X/K$, and an almost one-to-one
factor map $\rho:\widetilde Y\to Y$.

By the uniform Ellis group lift, $K$ acts strongly continuously on a cyclic
Hilbert realization of $X$ by unitaries commuting with $U$.  The tensor-moment
quotient lemma therefore shows that $\widetilde Y=X/K$ is
Hilbert-representable.  The already constructed map
$\rho:\widetilde Y\to Y$ is almost one-to-one, which is exactly option (c).
\end{proof}

\section*{Verification, limitations, and novelty}

The proof uses no selection theorem and no computational assumption.  The
moment map is defined on the whole quotient, not merely on a dense orbit.  Its
separation argument uses all tensor degrees, so it does not assume finite
generation of compact-group invariants.

The conclusion does not establish that $Y$ itself is Hilbert-representable
(option (a)), that $Y$ is a compact group-factor of a Hilbert-representable
system (option (b)), or the algebraic assertion in Problem 8.2.

Bounded searches on 9 August 2026 used the exact source question and arXiv id,
``factor of a Hilbert-representable system'', ``almost 1-1'' with Hilbert
representability, compact-group quotients, and symmetric tensor moments.  They
found later surveys repeating the question.  Ben Yaacov--Ibarluc\'ia--Tsankov
\cite{BIT} prove factor closure for the Hilbert compactification of a special
pro-oligomorphic class, not the general quotient lemma or option (c) above.  No
later general answer was found.  This is a bounded search, not a definitive
priority claim.

\section*{Human-review recommendation}

Review first the lift of $\Lambda_X$ to the strong unitary closure, especially
the weak-plus-equal-norm step on the cyclic transitive vector.  Then check that
equality of the averaged tensor moments determines the Haar orbit measure via
real Stone--Weierstrass.  These are the only non-source ingredients carrying
the conclusion.

\begin{thebibliography}{9}
\bibitem{GW}
E. Glasner and B. Weiss,
\emph{On Hilbert dynamical systems}, arXiv:1005.0230, 2010.

\bibitem{BIT}
I. Ben Yaacov, T. Ibarluc\'ia, and T. Tsankov,
\emph{Eberlein oligomorphic groups},
Trans. Amer. Math. Soc. 370 (2018), 2181--2209; arXiv:1602.05097.

\bibitem{Rudin}
W. Rudin, \emph{Real and Complex Analysis}, 3rd ed.,
McGraw--Hill, 1987.
\end{thebibliography}

\end{document}
